\section{Part D. Boyer--Moore Family}
\begin{frame}{Boyer--Moore Family의 전환}
\smission{pattern의 오른쪽부터 비교하고 mismatch 정보를 이용해 여러 text position을 건너뛴다.}
\begin{itemize}\item pattern은 text에 align된다.\item comparison은 \(j=m-1\)에서 왼쪽으로 진행한다.\end{itemize}
\begin{center}
\small
\textbf{pattern}
\(\rightarrow\)
\textbf{shift table}
\(\rightarrow\)
\textcolor{AlgoOrange}{\textbf{mismatch}}
\(\rightarrow\)
\textbf{shift distance}
\(\rightarrow\)
\textbf{next alignment}
\end{center}
\end{frame}
\begin{frame}{오른쪽부터 비교}
\centering\begin{tikzpicture}[x=.86cm,y=.78cm]
\only<1|handout:0>{%
  \smrow{2.1}{T}{C,A,T,I,G}
  \smshiftrow{.8}{P}{0}{T,I,G,E,R}
  \node[smbad]at(4,2.1){G};\node[smbad]at(4,.8){R};
  \draw[->,very thick,AlgoOrange](4,1.35)--(4,1.72);
}
\only<2|handout:1>{%
  \smrow{2.1}{T}{X,T,I,G,E,R}
  \smshiftrow{.8}{P}{1}{T,I,G,E,R}
  \foreach \x/\c in {1/T,2/I,3/G,4/E,5/R}{\node[smverified]at(\x,2.1){\c};\node[smverified]at(\x,.8){\c};}
  \draw[->,very thick,AlgoOrange](5,1.35)--(5,1.72);
  \node[font=\small,text=AlgoOrange]at(3,-.05){비교 방향: right \(\rightarrow\) left};
}
\node[smnote,text width=.76\linewidth]at(2.1,-1.0){%
\only<1|handout:0>{끝 문자 mismatch만으로도 큰 shift의 근거를 얻을 수 있다.}%
\only<2|handout:1>{오른쪽 suffix부터 확인하며, mismatch가 나면 문자와 suffix 정보를 사용한다.}};
\end{tikzpicture}
\end{frame}
\begin{frame}{Full Boyer--Moore vs Horspool}
\small\begin{tabular}{lll}\toprule
&Full BM&Horspool\\\midrule
bad-character&사용&단순화해 사용\\
good-suffix&사용&사용하지 않음\\
shift 기준&mismatch character/window&window 마지막 character\\
구현&복잡&간결\\\bottomrule
\end{tabular}
\end{frame}
\begin{frame}{Boyer--Moore Family Checkpoint}
\begin{enumerate}\item 비교 방향은?\item Horspool이 사용하지 않는 table은?\item 둘은 같은 알고리즘인가?\end{enumerate}
\begin{exampleblock}{답}오른쪽에서 왼쪽; good-suffix; 같은 family지만 동일하지 않다.\end{exampleblock}
\end{frame}
